\section{Discussion}
\label{sec:discussion}

\textbf{What the combination contributes.} Our starting components are not new individually. What is new is their joint behaviour on a graph-level visual operator. A linear pooling matrix leaves the error identity exact, so a formulation derived for node classification transfers without approximation; but a component that is beneficial in isolation can become harmful in the combination. The closed-form bias $\beta^\star$ cancels floor-quantisation error for an integrate-and-fire rate, whereas the deterministic two-phase schedule already delivers the exact count. Keeping both double-corrects the current. Identifying and removing this redundancy is a mechanistic contribution, not a hyperparameter choice: it is predicted by~\eqref{eq:dtp}, verified by a paired ablation, and consistent across both datasets and all seeds.

\textbf{Why the energy result is stable.} The 45\,nm per-operation ratio is only \MACACRatio, so a $5\times$ advantage requires the SNN's total ACs to fall below the ANN's MACs. Two mechanisms achieve this together: a narrow trained network reduces the number of emitting units, and $T\in\{2,3,4\}$ bounds the number of propagation steps. Neither alone suffices. Reporting a conservative convention that also charges real-valued propagation as MACs is important, because an AC-only count can be made arbitrarily favourable by driving the network toward zero spikes. We explicitly exclude zero-spike and chance-level-teacher configurations; those were the two ways our own screening initially produced false positives.

\textbf{The accuracy--energy trade-off is the result.} We do not claim a 1.5\,pp fidelity target. Under the $5\times$ constraint, measured conversion loss is \CIFARLoss$\pm$\CIFARLossStd\,pp on CIFAR10 and \MNISTLoss$\pm$\MNISTLossStd\,pp on MNIST, and teacher accuracy is \CIFARAnn\%/\MNISTAnn\%. A wider network or longer timestep improves fidelity but loses the energy advantage; a narrower one improves energy but loses fidelity. Increasing the training budget, which we expected to stabilise the teacher, instead widened margins that the low-$T$ quantiser amplifies and made fidelity worse (\EIGHTEENMNISTLoss/\EIGHTEENCIFARLoss\,pp). We report this negative result and the full Pareto frontier rather than choosing a favourable operating point after the fact.

\textbf{Limitations and threats to validity.} All energy figures are theoretical op-count estimates under a 45\,nm table; they ignore membrane-update cost, memory traffic and event-routing overhead, and are not measurements on neuromorphic hardware. Spike counts are therefore proxies. The method assumes a fixed non-negative linear adjacency and linear pooling; dynamic $k$-NN graphs, per-sample attention and max-relative aggregation break the factorisation in~\eqref{eq:identity} and are failure boundaries. Learning per-edge attention would require freezing it into a fixed operator first. Our teachers are deliberately small, so absolute accuracies are below published superpixel-GNN results; our claim concerns conversion fidelity and energy, not classifier quality. Finally, the four holdout groups are disjoint seeds drawn from the same two public benchmarks, so the conclusion is bounded to SuperPixel MNIST/CIFAR10 and does not establish ImageNet-scale behaviour. CUDA sparse reduction is non-deterministic; we mitigate it with five seeds, standard deviations and confidence intervals, and never select on a test split.
